%%
%% spring-lecture08.tex
%% 
%% Made by Alex Nelson <pqnelson@gmail.com>
%% Login   <alex@lisp>
%% 
%% Started on  2026-04-28T08:28:38-0700
%% Last update 2026-04-28T08:28:38-0700
%% 

\lecture{}

\begin{proposition}\label{prop:3.12}
Let $A$ be a simple algebra. The following are equivalent:
\begin{enumerate}
\item $A$ is semisimple
\item $A$ is right Artinian (i.e., $A$ is a right Artinian $A$-module)
\item $A$ has a right minimal ideal.
\end{enumerate}
\end{proposition}

\begin{proof}
$(1)\implies(2)$ is Definition~\ref{defn:3.1} (which implies $A$ is
  finitely-generated), then Proposition~\ref{prop:2.29} gives the
  result.

$(2)\implies(3)$ is the  definition of Artinian.

$(3)\implies(1)$ Let $N$ be a right minimal ideal of $A$. Then
\begin{equation}
AN=\sum_{x\in A}xN
\end{equation}
is a nonzero two-sided ideal of $A$. This must mean
\begin{equation}
A=AN.
\end{equation}
By Lemma~\ref{lemma:3.6}, for each $x\in A$ we have $xN$ is a simple
$A$-module, and then $A$ is semisimple by Proposition~\ref{prop:2.15}.
\end{proof}

\begin{proposition}\label{prop:3.13}
Let $A$ be a semisimple algebra. Then the following are equivalent:
\begin{enumerate}
\item $A$ is a simple algebra
\item All right minimal ideals of $A$ are isomorphic to each other
\item All simple $A$-modules are isomorphic to each other.
\end{enumerate}
\end{proposition}

\begin{proof}
$(2)\iff(3)$ by Proposition~\ref{prop:3.2}~(2).

$(1)\implies(2)$
For any minimal right ideals $N_{1}$, $N_{2}$ of $A$ we have
$AN_{1}=AN_{2}=A$ by the discussion from the previous Proposition. In particular,
\begin{subequations}
  \begin{align}
N_{1}N_{2} &= (N_{1}A)N_{2}\\
&=N_{1}(AN_{2})\\
&=N_{1}A\\
&=N_{1}\neq0,
  \end{align}
\end{subequations}
and thus (from last lecture) we have $N_{1}\iso N_{2}$ by Proposition~\ref{prop:3.8}.

$(2)\implies(1)$
Let $J$ be the nonzero two-sided ideal of $A$, so $0\neq J\in I(A)$.
Since $A$ is semisimple, there exists a minimal right ideal $N$ of $A$
such that $N\subset J$. By hypothesis (2) and using
Proposition~\ref{prop:3.8} and Proposition~\ref{prop:2.15} shows that
\begin{equation}
A=\sum_{x\in A}xN\subset J,
\end{equation}
which means $J=A$. Hence $A$ is a simple algebra.
\end{proof}

\begin{corollary}\label{cor:3.14}
Let $A$ be a semisimple algebra. Let $N$ be a minimal right ideal of
$A$. For any $A$-module $M$, there exists a unique cardinal number
$\alpha$ such that $M\iso\bigoplus\alpha N$ is the direct sum of
$\alpha$ copies of $N$.
\end{corollary}

\begin{proof}
By Propositions~\ref{prop:3.12} for semisimple,~\ref{prop:3.13} for
simple and isomorphic, and Proposition~\ref{prop:3.2}~(2),
and Proposition~\ref{prop:2.15}.
\end{proof}

\begin{corollary}\label{cor:3.15}
Let $A$ be a finite-dimensional simple $\FF$-algebra (where $\FF$ is
some field). Let $M_{1}$, $M_{2}$ be $A$-modules. Then $M_{1}\iso M_{2}$
as $A$-modules if and only if $\dim_{\FF}(M_{1})=\dim_{\FF}(M_{2})$.
\end{corollary}

\begin{proof}
Since $A$ is finite-dimensional, $A$ is right Artinian by
Proposition~\ref{prop:3.12} and has a minimal right ideal $N$ (such
that $\dim_{\FF}N$ is finite). By Corollary~\ref{cor:3.14},
$M_{1}\iso\bigoplus\alpha N$ and $M_{2}\iso\bigoplus\beta N$ so
$M_{1}\iso M_{2}$ if and only if $\alpha=\beta$ if and only if
\begin{equation}
\dim_{\FF}M_{1}=\alpha\dim_{\FF}N=\beta\dim_{\FF}N=\dim_{\FF}M_{2}.
\end{equation}
Hence the claim.
\end{proof}

\begin{example}\label{ex:3.16}
For a division algebra $D$, let $A=M_{n}(D)$ be a matrix algebra, let
$N_{i}=\canMat_{i,i}A$ for $i=1,\dots,n$ and the canonical matrix
$\canMat_{i,i}$ (with its only nonzero entry being 1 at row $i$
and column $i$).
\begin{enumerate}
\item These $N_{i}$ are minimal right ideals of $A$
\item $A=\bigoplus^{n}_{i}N_{i}$
\item For all $i,j$ we have $N_{i}\iso N_{j}$
\item $A$ is simple and semisimple
\item $\End_{A}(N_{i})\iso D$
\end{enumerate}
\end{example}

\begin{proof}
Homework.
\end{proof}

\begin{notation}
Let $A$ be an $R$-algebra. Let $M_{1}$, \dots, $M_{n}$ be $A$-modules.
We denote by $[\hom_{A}(M_{j},M_{i})]$ for the matrix
\begin{equation}
[\hom_{A}(M_{j},M_{i})]:=\begin{bmatrix}\hom_{A}(M_{1},M_{1}) & \dots
& \hom_{A}(M_{1},M_{n}\\
\vdots & \ddots & \vdots\\
\hom_{A}(M_{n},M_{1}) & \dots & \hom_{A}(M_{n},M_{n})
\end{bmatrix}
\end{equation}
Further, if $[\varphi_{i,j}]\in[\hom_{A}(M_{j},M_{i})]$, we see that
$[\hom_{A}(M_{j},M_{i})]$ forms an $R$-algebra by
$[\phi_{i,j}][\psi_{j,k}]=[\chi_{i,k}]$ where
\begin{equation}
\chi_{i,k}=\sum^{n}_{j=1}\phi_{i,j}\circ\psi_{j,k}
\end{equation}
is ``just'' matrix multiplication.
\end{notation}

\begin{proposition}\label{prop:3.17}
$[\hom_{A}(M_{j},M_{i})]\iso\End_{A}(\bigoplus^{n}_{i=1}M_{i})$ as $R$-algebras.
\end{proposition}

\begin{proof}
Let $M=\bigoplus^{n}_{i=1}M_{i}$, let $\pi_{j}\colon M\to M_{j}$ be
the projection morphism for $j=1,\dots,n$. Let $\kappa_{j}\colon M_{j}\into M$
be the inclusion morphism. Then we have
\begin{equation}
\sum^{n}_{j=1}\kappa_{j}\circ\pi_{j}=\id_{M}
\end{equation}
and
\begin{equation}
\pi_{i}\circ\kappa_{j}=\delta_{ij}\id_{M_{j}}
\end{equation}
(so $\pi_{i}\circ\kappa_{j}$ is the zero map for $i\neq j$). Let us
define the mappings
\begin{equation}
\begin{split}
\alpha\colon&\End_{A}(M)\to[\hom_{A}(M_{j},M_{i})]\\
&\phi\mapsto[\pi_{i}\circ\kappa_{j}\circ\phi]
\end{split}
\end{equation}
and
\begin{equation}
\begin{split}
\beta\colon&[\hom_{A}(M_{j},M_{i})]\to\End_{A}(M)\\
&[\phi_{i,j}]\mapsto\sum^{n}_{i,j=1}\kappa_{i}\circ\phi_{i,j}\circ\pi_{j}.
\end{split}
\end{equation}
Then by direct calculation, we can verify
\begin{equation}
\alpha\circ\beta=\id\quad\mbox{and}\quad\beta\circ\alpha=\id.
\end{equation}
Further, $\alpha$ is a morphism since
\begin{subequations}
  \begin{align}
\alpha(\phi\psi) &= [\pi_{i}\phi\psi\kappa_{k}] = [\pi_{i}(\sum_{j}\kappa_{j}\pi_{j})\psi\kappa_{k}]\\
&=[\sum_{j}(\pi_{i}\phi\kappa_{j})(\pi_{j}\psi\kappa_{k})]\\
&=\alpha(\phi)\alpha(\psi)
  \end{align}
\end{subequations}
by matrix multiplication.
\end{proof}

\begin{corollary}\label{cor:3.18}
Let $A$ be an $R$-algebra. Let $M_{1}$, \dots, $M_{n}$, $M$ be
$A$-modules. Then:
\begin{enumerate}
\item $\End_{A}(\bigoplus nM)\iso M_{n}(\End_{A}(M))$
\item If $M$ is a free $A$-module with $n$ generators, then
  $\End_{A}(M)\iso M_{n}(A)$
\item If $\hom_{A}(M_{j},M_{i})=0$ for all $i\neq j$, then
  $\End_{A}(\bigoplus^{n}_{i=1}M_{i})\iso\prod^{n}_{i=1}\End_{A}(M_{i})$
\end{enumerate}
\end{corollary}

\begin{proof}
(1) and (3) are obvious.

(2) Since $M$ is free, $M\iso\bigoplus nA$ and the result follows by Definition~\ref{defn:1.6}.
\end{proof}

\begin{theorem}[Wedderburn's structure theorem]\label{thm:3.19}
Let $A$ be a right or left semisimple $R$-algebra. Then the following
all hold:
\begin{enumerate}
\item There exists natural numbers $n_{1},\dots,n_{r}\in\NN$ and
  $R$-division algebras
  $D_{1}$, \dots, $D_{r}$ such that $A\iso M_{n_{1}}(D_{1})\times\cdots\times M_{n_{r}}(D_{r})$
\item The pairs $(n_{i},D_{i})$ are uniquely determined (up to
  isomorphism) by $A$
\item Conversely for any $n_{1}$, \dots, $n_{r}\in\NN$ and for any
  division algebras (over $R$) $D_{1}$, \dots, $D_{r}$, we have
  $M_{n_{1}}(D_{1})\times\cdots\times M_{n_{r}}(D_{r})$ is a right
  \emph{and} left $R$-algebra (i.e., the class of right
  semisimple algebras and left semisimple algebras coincide).
\end{enumerate}
\end{theorem}

\begin{proof}
(1) $A\iso\bigoplus^{r}_{i=1}M_{i}$ where ach $M_{i}\iso\bigoplus n_{i}N_{i}$
and $N_{i}$ are minimal right ideals of $A$ such that $N_{i}\not\iso N_{j}$
for $i\neq j$. By Lemma~\ref{lemma:2.18}~(2) [if $\phi\colon A\to B$ and
  $A=\bigoplus M_{i}$ and $B\iso\bigoplus M'_{i}$, then
  $\phi(M_{i})\subset M'_{i}$], we have $\hom_{A}(M_{i},M_{j})=0$ for
$i\neq j$. The isomorphism follows from
\begin{subequations}
  \begin{align}
A&\iso\End_{A}(A)\quad\mbox{by Def~\ref{defn:1.6}}\\
&\iso\prod^{r}_{i=1}\End_{A}(M_{i})\quad\mbox{by Cor~\ref{cor:3.18}~(3)}\\
&\iso\prod^{r}_{i=1}M_{n_{i}}(D_{i})\quad\mbox{by Cor~\ref{cor:3.18}~(1) where }
D_{i}=\End_{A}(N_{i}).
  \end{align}
\end{subequations}
We see $D_{i}$ is a division algebra by
Corollary~\ref{cor:2.12}. Hence the claim for the right case. The left
case is similar.

(2) Suppose $A=A_{1}\times\cdots\times A_{s}$ where each $M_{i}\iso M_{k_{i}}(C_{i})$
is a matrix algebra for some $k_{i}\in\NN$ and $C_{i}$ is a division
algebra over $R$. By Example~\ref{ex:3.16} and Lemma~\ref{lemma:3.4}
$A_{i}$ is semisimple so $A_{i}\iso\bigoplus k_{i}P_{i}$ as
$A$-modules, where $P_{i}$ is a minimal right ideal such that
$P_{i}\subset A_{i}$ and
\begin{equation}
C_{i}\iso\End_{A_{i}}(P_{i})\iso\End_{A}(P_{i}).
\end{equation}
Since $P_{i}P_{j}\subset A_{i}A_{j}=0$ (since
$A_{i}=0\times\cdots\times0\times A_{i}\times0\times\cdots\times0$ and
$A_{j}=0\times\cdots\times0\times A_{j}\times0\times\cdots\times0$,
so their product $A_{i}A_{j}=0$ if $i\neq j$) by Proposition~\ref{prop:3.8},
$P_{i}\not\iso P_{j}$ if $i\neq j$.
By applying Proposition~\ref{prop:2.19} to
\begin{equation}
A=\bigoplus^{r}_{i=1}\left(\bigoplus
n_{i}N_{i}\right)=\bigoplus^{s}_{i=1}\left(\bigoplus k_{i}P_{i}\right),
\end{equation}
we must have $(n_{i},N_{i})\iso(k_{j},P_{j})$ so
\begin{equation}
D_{i}=\End_{A}(N_{i})\iso\End_{A}(P_{j})=C_{j},
\end{equation}
proving the claim up to permutation.

(3) By Example~\ref{ex:3.16} and its left version, $M_{n_{i}}(D_{i})$
are left and right semisimple $R$-algebras, and thus $\prod^{r}_{i=1}M_{n_{i}}(D_{i})$
is a left and right algebra by Corollary~\ref{cor:3.5}.
\end{proof}

\begin{corollary}\label{cor:3.20}
We have $A$ be a right or left Artinian simple algebra if and only if
there exists a unique $n\in\NN$ and a unique division algebra $D$ such
that $A\iso M_{n}(D)$.
\end{corollary}

\begin{proof}
By Proposition~\ref{prop:3.12}, Proposition~\ref{prop:3.13}, and Theorem~\ref{thm:3.19}.
\end{proof}